\documentclass[../../main.tex]{subfiles}
\begin{document}

\begin{center}
  {\fontsize{48pt}{58pt}\selectfont \textbf{2007}\normalsize\textbf{年度}}
\end{center}

\begin{table}[htb]
  \centering
  \caption{出題テーマと難易度}
  \begin{tabular}{|c|p{9cm}|c|}
    \hline
    \textbf{問題番号} & \textbf{テーマ} & \textbf{難易度} \\
    \hline
    第1問 & 整式の係数の整数性 & \\
    \hline
    第2問 & 平面図形の辺の長さと極限 & \\
    \hline
    第3問 & 放物線上の二点を内分する点の動く範囲 & \\
    \hline
    第4問 & 行列の分解と行列の積 & \\
    \hline
    第5問 & 硬貨投げとブロックの高さの確率 & \\
    \hline
    第6問 & 定積分の評価と$\log 2$の不等式 & \\
    \hline
  \end{tabular}
\end{table}

\begin{table}[htb]
  \centering
  \caption{解答}
  \renewcommand{\arraystretch}{1.4}
  \begin{tabular}{|c|c|p{9cm}|}
    \hline
    \textbf{問題番号} & \textbf{小問} & \textbf{答え} \\
    \hline
    第1問 & -- & $P(x)$の$n$次以下の係数は整数（証明問題） \\
    \hline
    第2問 & -- & $\displaystyle\lim_{n\to\infty}s_n
    =(e-1)\sqrt{\pi^2+1}$ \\
    \hline
    第3問 & (1) & $-1\le a\le1$，$a^2\le b\le B(a)$．ただし
    $B(a)=
    \begin{dcases}
      \dfrac{3}{2}a^2+|a|+\dfrac{1}{2} & (|a|\le\dfrac{1}{3}) \\
      3a^2-4|a|+2 & (\dfrac{1}{3}\le|a|\le1)
    \end{dcases}$ \\
    \cline{2-3}
    & (2) & 領域$D$を図示（解答参照） \\
    \hline
    第4問 & (1)〜(3) & \\
    \hline
    第5問 & (1) & $m<n$なら$p_m=(1-p)p^m$，$m=n$なら$p_m=p^n$ \\
    \cline{2-3}
    & (2) & $m<n$なら$q_m=1-p^{m+1}$，$m=n$なら$q_m=1$ \\
    \cline{2-3}
    & (3) & $m<n$なら$r_m=(1-p)p^m(2-p^m-p^{m+1})$，
    $m=n$なら$r_m=p^n(2-p^n)$ \\
    \hline
    第6問 & (1) & 定積分の上下界を示す（証明問題） \\
    \cline{2-3}
    & (2) & $0.68<\log 2<0.71$（証明問題） \\
    \hline
  \end{tabular}
\end{table}

\paragraph{試験形態}
6問．試験時間は要確認．

\end{document}
